\documentclass[11pt,aspectratio=169]{beamer}
\ifdefined\pdfmapfile\pdfmapfile{+cm-super-t1.map}\fi
% Shared preamble for MA3521 lectures.
\usepackage{beamerthemeMTU}
\usepackage{amsmath,amssymb,cancel}
\course{MA3521 --- Introduction to ODEs}
\author{Allan Struthers}
\institute{Math Sciences}
% Reusable shear-building sketch. Displacements are absolute, measured from rest.
% Numeric data: masses in kg, stiffnesses in N/m, damping in N s/m.
% Drawing offsets are independent of physical units.
% basemotion accepts a math expression, for example basemotion={\sin(t)}.
\usetikzlibrary{arrows.meta,calc}
\pgfkeys{/ode building/.is family,/ode building,
  baseoffset/.store in=\ODEbase, basemotion/.store in=\ODEbasemotion,
  offset1/.store in=\ODEbone,offset2/.store in=\ODEbtwo,
  mass1/.store in=\ODEbmone,mass2/.store in=\ODEbmtwo,
  stiffness1/.store in=\ODEbkone,stiffness2/.store in=\ODEbktwo,
  damping1/.store in=\ODEbcone,damping2/.store in=\ODEbctwo}

% Dashpot between two numeric points. The piston follows the line of action.
\newcommand{\ODEBuildingDamper}[4]{%
  \pgfmathsetmacro{\ODEangle}{atan2(#4-#2,#3-#1)}
  \pgfmathsetmacro{\ODElength}{sqrt((#3-#1)^2+(#4-#2)^2)}
  \begin{scope}[shift={(#1,#2)},rotate=\ODEangle]
    \draw[thick] (0,0)--({.38*\ODElength},0);
    \draw[thick,fill=white] ({.38*\ODElength},-.075) rectangle ({.72*\ODElength},.075);
    \draw[thick] ({.55*\ODElength},-.065)--({.55*\ODElength},.065);
    \draw[thick] ({.55*\ODElength},0)--(\ODElength,0);
  \end{scope}}

\newcommand{\TwoStoryBuilding}[1]{%
  \begingroup
  \pgfkeys{/ode building/.cd,
    basemotion={\sin(t)},baseoffset=.30,offset1=.65,offset2=1.05,
    mass1=4,mass2=3,
    stiffness1=4,stiffness2=3,
    damping1=2,damping2=1,#1}
  \begin{tikzpicture}[font=\normalsize,
    motion/.style={-{Stealth[length=2mm]},thick,mtuGreen},
    ref/.style={gray!55,dashed,thin},
    tag/.style={fill=white,inner sep=1.5pt}]
    % Displacement arrows extend from the outside datum to the left floor edges.
    % The datum is offset for legibility. Arrows are schematic, not to scale.
    % Shape offsets control the cubic column geometry independently.
    \draw[ref] (-3.15,-.05)--(-3.15,3.02);
    \node[anchor=south,font=\small] at (1.5,3.38) {Displacements: $\vect{y}(t)=(y_1(t),y_2(t))$ (m),\quad $t$ in seconds.};
    \node[gray,anchor=north,font=\small] at (-3.15,-.1) {$0$};
    % Light rest-position outline.

    \foreach \i in {1,2}{\draw[ref] (-1.2,{1.35*\i})--(1.2,{1.35*\i});}
    % Slabs end flush with the columns. No overhanging edges.
    \foreach \i/\offset in {1/\ODEbone,2/\ODEbtwo}{
      \draw[thick,fill=gray!25] ({\offset-1.2},{1.35*\i-.055}) rectangle ({\offset+1.2},{1.35*\i+.055});


      \draw[motion] (-3.15,{1.35*\i})--({\offset-1.2},{1.35*\i});
      \node[above,mtuGreen] at (-2.6,{1.35*\i}) {$y_{\i}(t)$};

    }
    % Each outside edge is ONE continuous piecewise cubic path. Vertical
    % tangents and matching derivatives at the floor junctions give C1 joins.
    \foreach \side in {-1.2,1.2}{
      \draw[line width=1.7pt] ({\ODEbase+\side},0)
      .. controls ({\ODEbase+\side},.45) and ({\ODEbone+\side},.90) .. ({\ODEbone+\side},1.35)
      .. controls ({\ODEbone+\side},1.80) and ({\ODEbtwo+\side},2.25) .. ({\ODEbtwo+\side},2.70);
    }
    \foreach \i/\lo/\hi/\stiff/\damp/\mass in {
      1/\ODEbase/\ODEbone/\ODEbkone/\ODEbcone/\ODEbmone,
      2/\ODEbone/\ODEbtwo/\ODEbktwo/\ODEbctwo/\ODEbmtwo}{
      \pgfmathsetmacro{\ODElowy}{1.35*(\i-1)}
      \pgfmathsetmacro{\ODEhighy}{1.35*\i}
      \pgfmathsetmacro{\ODEmid}{(\lo+\hi)/2}
      \node[tag] at ({\ODEmid-.87},{\ODElowy+.67}) {$k_{\i}$};
      \ODEBuildingDamper{\lo+.2}{\ODElowy+.055}{\hi+.65}{\ODEhighy-.055}
      \node[anchor=west,tag] at ({\ODEmid+.52},{\ODElowy+.66}) {$c_{\i}$};
      \node[anchor=west,font=\small,draw=mtuGreen,line width=.6pt,inner xsep=4pt,inner ysep=1pt,minimum width=2.7cm] at (3.35,{\ODEhighy-.32+.2*(\i-2)}) {\setlength{\jot}{0pt}$\begin{aligned}
        m_{\i}&=\mass\,\mathrm{kg}\\k_{\i}&=\stiff\,\mathrm{N/m}\\c_{\i}&=\damp\,\mathrm{N\,s/m}
      \end{aligned}$};
    }
    % Foundation and rollers. The baseoffset uses the same outside strip.
    \draw[thick,fill=gray!40] ({\ODEbase-1.5},-.16) rectangle ({\ODEbase+1.5},0);
    \foreach \wheel in {-.95,.95}{\draw[thick,fill=white] ({\ODEbase+\wheel},-.31) circle[radius=.15];}
    \draw[thick] (-1.8,-.47)--(2.4,-.47);
    \foreach \ground in {-1.7,-1.45,...,2.3}{\draw[gray,thin] (\ground,-.47)--({\ground-.12},-.6);}

    \draw[motion] (-3.15,-.02)--({\ODEbase-1.5},-.02);
    \node[above,mtuGreen] at (-2.6,-.02) {$y_0(t)$};
    \node[anchor=west,align=left,font=\small] at (3.35,-.13) {$y_0(t)=\ODEbasemotion\;\mathrm{m}$};
  \end{tikzpicture}%
  \endgroup}

\title{Variation of Parameters and Two-Story Buildings}
\date{}
\lecturelabel{Review}

% Recap follows Week 10, 9_Lecture.tex.
% Building problems adapt the two cases in Week 9, BuildingWorksheet.tex.

\begin{document}
\begin{frame}{Variation of Parameters for a Scalar Equation}
\[
  y'+p(t)y=q(t)
  \qquad
  y_1'=-p(t)y_1
  \qquad
  y_c=cy_1
\]
\keypoint{Variation of parameters} replaces the constant $c$ by a function $u(t)$.
\[
  y=uy_1
  \qquad
  y'+p(t)y=u'y_1+u\underbrace{(y_1'+p(t)y_1)}_{0}
\]
So $u'=q(t)/y_1(t)$. Integrate and use the initial condition.

\begin{block}{Example \quad $y'+2y=e^{-2t}$,\quad $y(0)=1$}
Here $y_1(t)=e^{-2t}$, so $u'=1$ and $u(0)=1$.
Then $u(t)=1+t$ and $y(t)=(1+t)e^{-2t}$.
\end{block}
\end{frame}

\begin{frame}{Variation of Parameters for a System}
For $\vect{y}'=\mat{A}\vect{y}+\vect{f}(t)$, the
\keypoint{fundamental matrix} has $n$ linearly independent complementary solutions as its columns.
\[
  \mat{\Phi}(t)=[\,\vect{y}_1(t)\ \cdots\ \vect{y}_n(t)\,]
  \qquad
  \mat{\Phi}'=\mat{A}\mat{\Phi}
  \qquad
  \vect{y}_c=\mat{\Phi}\vect{c}
\]
Replace $\vect{c}$ by $\vect{u}(t)$ and substitute $\vect{y}=\mat{\Phi}\vect{u}$.
\[
  \underbrace{(\mat{\Phi}'-\mat{A}\mat{\Phi})}_{0}\vect{u}
  +\mat{\Phi}\vect{u}'=\vect{f}(t)
  \qquad
  \vect{u}'=\mat{\Phi}^{-1}\vect{f}(t)
\]
At each $t$, the last equation amounts to solving
$\mat{\Phi}(t)\vect{u}'(t)=\vect{f}(t)$.
\end{frame}

\begin{frame}{Initial Conditions and the Forcing Integral}
Given $\vect{y}(0)=\vect{y}_0$, the coefficient vector starts at
$\vect{u}(0)=\mat{\Phi}(0)^{-1}\vect{y}_0$.
\[
  \boxed{\vect{y}(t)=
    \underbrace{\mat{\Phi}(t)\mat{\Phi}(0)^{-1}\vect{y}_0}_{\text{response to initial data}}
    +\underbrace{\mat{\Phi}(t)\int_0^t\mat{\Phi}(s)^{-1}\vect{f}(s)\,ds}_{\text{response to forcing}}}
\]
The first term makes $\vect{y}(0)=\vect{y}_0$ since the integral is zero at $t=0$. 
For $\mat{A}$ constant
\[
  \mat{\Phi}(t) = e^{\mat{A}t}=\sum_{k=0}^{\infty}\frac{(\mat{A}t)^k}{k!}
    =\mat{I}+\mat{A}t+\frac{(\mat{A}t)^2}{2!}+\cdots
\]
and the solution of the IVP is
\[
  \vect{y}(t)=e^{\mat{A}t}\vect{y}_0
    +\int_0^t e^{\mat{A}(t-s)}\vect{f}(s)\,ds
\]
\end{frame}

\begin{frame}{Two-Story Building}
\begin{columns}[c,onlytextwidth]
  \column{.61\textwidth}
  \resizebox{\linewidth}{!}{\TwoStoryBuilding{}}
  \column{.35\textwidth}
  \textbf{Worked problem}

  \medskip
  The base motion $y_0(t)$ is prescribed. The two floor displacements are
  measured to the right from rest.

  \[
    \vect{y}(0)=(0,-1)\,\mathrm{m}
  \]
  \[
    \vect{y}'(0)=(0,2)\,\mathrm{m/s}
  \]
\end{columns}
\end{frame}

\begin{frame}{Story Forces and Floor Balances}
\setlength{\abovedisplayskip}{5pt}
\setlength{\belowdisplayskip}{5pt}
\setlength{\abovedisplayshortskip}{5pt}
\setlength{\belowdisplayshortskip}{5pt}
Story $i$ connects levels $i-1$ and $i$. The force on the 
upper level is
$-f_i$ and the force on the lower level is $+f_i$, where
\[
  f_i=k_i(y_i-y_{i-1})+c_i(y_i'-y_{i-1}')
\]
\[
  \underbrace{\frac{\mathrm{N}}{\cancel{\mathrm{m}}}
    \cancel{\mathrm{m}}}_{\text{spring}}=\mathrm{N}
  \qquad
  \underbrace{\frac{\mathrm{N}\,\cancel{\mathrm{s}}}{\cancel{\mathrm{m}}}
    \frac{\cancel{\mathrm{m}}}{\cancel{\mathrm{s}}}}_{\text{damper}}=\mathrm{N}
  \qquad
  \mathrm{kg}\,\frac{\mathrm{m}}{\mathrm{s}^2}=\mathrm{N}
\]
Substituting the story forces into Newton's law gives
\[
  \begin{alignedat}{2}
    4y_1'' &=-f_1+f_2 &\quad
      &=-4(y_1-y_0)-2(y_1'-y_0')+3(y_2-y_1)+(y_2'-y_1')\\
    3y_2'' &=-f_2 &
      &=\phantom{-4(y_1-y_0)-2(y_1'-y_0')}-3(y_2-y_1)-(y_2'-y_1')
  \end{alignedat}
\]
\end{frame}

\begin{frame}{Floor Balances in Matrix Form}
Collect the floor terms on the left and the prescribed base motion on the right.
\[
  \underbrace{\begin{pmatrix}4&0\\0&3\end{pmatrix}}_{\mat{M}}
    \begin{pmatrix}y_1''\\y_2''\end{pmatrix}
  +\underbrace{\begin{pmatrix}3&-1\\-1&1\end{pmatrix}}_{\mat{C}}
    \begin{pmatrix}y_1'\\y_2'\end{pmatrix}
  +\underbrace{\begin{pmatrix}7&-3\\-3&3\end{pmatrix}}_{\mat{K}}
    \begin{pmatrix}y_1\\y_2\end{pmatrix}
  =\underbrace{\begin{pmatrix}4y_0+2y_0'\\0\end{pmatrix}}_{\vect{b}(t)}
\]
\[
  \mat{M}\vect{y}''+\mat{C}\vect{y}'+\mat{K}\vect{y}=\vect{b}(t)
  \qquad
  \vect{b}(t)=\begin{pmatrix}4y_0(t)+2y_0'(t)\\0\end{pmatrix}
\]
The units of $\mat{M}$, $\mat{C}$, $\mat{K}$, and $\vect{b}$ are
$\mathrm{kg}$, $\mathrm{N\,s/m}$, $\mathrm{N/m}$, and $\mathrm{N}$, respectively.

\medskip
For this example, the prescribed base motion is $y_0(t)=\sin(t)$.
\end{frame}

\begin{frame}{First-Order Building IVP}
Let $\vect{z}=(y_1,y_2,y_1',y_2')$.
The top two equations give velocity. The bottom two give acceleration.
\[
  \vect{z}'=\mat{A}\vect{z}+\vect{g}(t)
  \qquad
  \mat{A}=\begin{pmatrix}\mat{0}&\mat{I}\\-\mat{M}^{-1}\mat{K}&-\mat{M}^{-1}\mat{C}\end{pmatrix}
  \qquad
  \vect{g}(t)=\begin{pmatrix}\vect{0}\\\mat{M}^{-1}\vect{b}(t)\end{pmatrix}
\]
\[
  \mat{A}=\begin{pmatrix}
    0&0&1&0\\
    0&0&0&1\\
    -7/4&3/4&-3/4&1/4\\
    1&-1&1/3&-1/3
  \end{pmatrix}
  \quad
  \vect{g}(t)=\begin{pmatrix}0\\0\\(4y_0(t)+2y_0'(t))/4\\0\end{pmatrix}
  \quad
  \vect{z}(0)=\begin{pmatrix}0\\-1\\0\\2\end{pmatrix}
\]
Dividing each force by its floor mass gives acceleration in
$\mathrm{N}/\mathrm{kg}=(\cancel{\mathrm{kg}}\,\mathrm{m}/\mathrm{s}^2)/\cancel{\mathrm{kg}}=\mathrm{m}/\mathrm{s}^2$.
\end{frame}

\begin{frame}{Two-Story Building Practice \quad $\mat{M}\vect{y}''+\mat{C}\vect{y}'+\mat{K}\vect{y}=\vect{b}(t)$}
\begin{columns}[c,onlytextwidth]
  \column{.59\textwidth}
  \resizebox{\linewidth}{!}{\TwoStoryBuilding{%
    basemotion={2\sin(t)},mass1=3,mass2=2,%
    stiffness1=6,stiffness2=4,damping1=3,damping2=2}}
  \column{.38\textwidth}
  \[
    \begin{aligned}
      \vect{y}(0)&=(0,1)\,\mathrm{m}\\
      \vect{y}'(0)&=(-1,0)\,\mathrm{m/s}
    \end{aligned}
  \]
  \begin{enumerate}
    \item Find $\mat{M}$, $\mat{C}$, $\mat{K}$, and $\vect{b}(t)$.
    \item Write $\vect{z}'=\mat{A}\vect{z}+\vect{g}(t)$ and $\vect{z}(0)$.
  \end{enumerate}
  Stop once the IVP is set up.
\end{columns}
\end{frame}

\begin{frame}{Practice Check \quad $\mat{M}\vect{y}''+\mat{C}\vect{y}'+\mat{K}\vect{y}=\vect{b}(t)$}
\[
  \mat{M}=\begin{pmatrix}3&0\\0&2\end{pmatrix}
  \quad
  \mat{C}=\begin{pmatrix}5&-2\\-2&2\end{pmatrix}
  \quad
  \mat{K}=\begin{pmatrix}10&-4\\-4&4\end{pmatrix}
  \quad
  \vect{b}(t)=\begin{pmatrix}12\sin(t)+6\cos(t)\\0\end{pmatrix}
\]
\[
  \mat{A}=\begin{pmatrix}
    0&0&1&0\\
    0&0&0&1\\
    -10/3&4/3&-5/3&2/3\\
    2&-2&1&-1
  \end{pmatrix}
  \quad
  \vect{g}(t)=\begin{pmatrix}0\\0\\4\sin(t)+2\cos(t)\\0\end{pmatrix}
  \quad
  \vect{z}(0)=\begin{pmatrix}0\\1\\-1\\0\end{pmatrix}
\]
As a check, the base term in the acceleration equation is
$\bigl(6y_0(t)+3y_0'(t)\bigr)/3$.
\end{frame}
\end{document}
